\begingroup
\clearpage
\let\clearpage\relax
\vspace*{-16pt}
\chapter[METHODOLOGIES]{METHODOLOGIES}\label{sec.methods}
\endgroup

This chapter develops six solution methods for the maximum mortality
rate clique problem, presented in order of increasing
sophistication. Section~\ref{sec.MBK} introduces a modified
Bron--Kerbosch enumerative algorithm that searches all $\ell$-frequent
cliques while tracking the best mortality rate found.
Section~\ref{sec.CBB} presents a combinatorial branch-and-bound
algorithm that augments the enumeration with an upper-bounding
scheme to enable tighter pruning. Section~\ref{sec.formulations}
develops two MILP formulations: one based on auxiliary product
variables and one based on the Charnes--Cooper transformation.
Section~\ref{sec.lbbd} introduces a logic-based Benders reformulation
that avoids the patient-indexed variables driving the scalability
limitations of the MILP formulations, and shows how it can be
strengthened through a hybrid strategy that integrates the
combinatorial branch-and-bound algorithm to generate valid
inequalities within a decomposition branch-and-cut framework.

\section{The Modified Bron--Kerbosch Enumerative
Algorithm}\label{sec.MBK}

The original~\cite{BronKerbosch73} (BK) algorithm is a recursive procedure
that operates on three sets: a current clique $C$; a candidate set
$L$ containing every vertex that may be added to enlarge $C$; and an
exclusion set $S$ containing vertices that are adjacent to all
vertices in $C$ but are excluded as candidates, in order to avoid
outputting the same maximal clique multiple times. The algorithm
selects a vertex from $L$, moves it into $C$, and updates both $L$
and $S$ by removing the non-neighbors of the added vertex $v$, before
invoking the recursive call with the updated sets. When both $L$ and
$S$ are empty, the current clique is maximal and is returned. The
algorithm then backtracks to the most recent recursive call, removes
the previously added vertex from $L$, and transfers it to $S$ to
prevent duplicate outputs.

Several variants for enumerating cliques have been introduced since
the classical algorithm of BK; see for
instance~\citep{chiba1985, TOMITA2006, CAZALS2008TCS, Eppstein13,
san2018efficiently}. Because the mortality rate function is not
monotone, we cannot restrict the search to maximal cliques, and many
enhancements to the classical algorithm, such as pivoting, do not
apply to MMRCP.  The algorithm introduced here follows the BK recursive structure and maintains the $\ell$-frequent clique with the highest mortality rate encountered during the search.

\begin{algorithm}[htbp]
    \SetKwFunction{FBronKerbosch}{Bron-Kerbosch}
    \SetKwProg{Fn}{Function}{}{}
    \KwIn{$\langle P, \widetilde{P}, A, G=(V,E), \ell\rangle$ and
    initial clique $C^0$ (possibly empty)}
    \KwOut{An $\ell$-frequent clique containing $C^0$ with maximum
    mortality rate}
    Initialize global variables $\texttt{BestSol} \gets C^0,\
    \texttt{BestObj} \gets \mu(C^0)$\\
    \If{$C^0 = \emptyset$}{
        $v \gets \arg\max_{u \in V} \mu(u)$\\
        $\texttt{BestSol} \gets \{v\},\ \texttt{BestObj} \gets
        \mu(v)$\\
        \FBronKerbosch{$\emptyset, V, P, \widetilde{P}$}
    }
    \Else{
        $L \gets \bigcap\limits_{u \in C^0} N(u) \setminus C^0$\\
        \FBronKerbosch{$C^0, L, \den(C^0), \num(C^0)$}
    }
    \KwRet{$\texttt{BestSol}$ and $\texttt{BestObj}$}\\
    \Fn{\FBronKerbosch{$C, L, \texttt{CurrentDen},
    \texttt{CurrentNum}$}}{
        \If{$L = \emptyset$ \label{step.Mterminate}}{\KwRet}
        \For{$v \in L$ \label{step.outer-for}}{
            $\texttt{NewDen} \gets \texttt{CurrentDen} \cap A_v$\\
            $\texttt{NewNum} \gets \texttt{CurrentNum} \cap A_v$\\
            \If{$|\texttt{NewDen}| \ge \ell$}{
                \If{$\frac{|\texttt{NewNum}|}{|\texttt{NewDen}|} >
                \texttt{BestObj}$}{
                    $\texttt{BestObj} \gets
                    \frac{|\texttt{NewNum}|}{|\texttt{NewDen}|}$\\
                    $\texttt{BestSol} \gets C \cup \{v\}$\\
                }
                \FBronKerbosch{$C \cup \{v\}, L \cap N(v),
                \texttt{NewDen}, \texttt{NewNum}$}\\
            }
            $L \gets L \setminus \{v\}$\\
        }
        \KwRet
    }
\caption{Modified {Bron--Kerbosch} for the Maximum
Mortality Rate Clique Problem}
\label{alg.BKMMRC}
\end{algorithm}

Algorithm~\ref{alg.BKMMRC} solves MMRCP~\eqref{eq.MMRCP} and includes
a pruning step based on the following feasibility condition: a branch
is discarded whenever $\left|\bigcap_{u \in C} A_u\right| < \ell$,
i.e., whenever the current partial clique is not $\ell$-frequent. The
notation $N(v) \coloneqq \{u : \{u,v\} \in E\}$ denotes the
neighborhood of vertex $v$ in $G$.

Global variables \texttt{BestSol} and \texttt{BestObj} track the
incumbent clique and its mortality rate throughout the search, and are
updated whenever a feasible clique with a strictly higher mortality
rate is encountered. If $C^0$ is nonempty, it initializes the
incumbent; otherwise, the single-vertex clique with the highest
mortality rate serves as the initial incumbent. The candidate set $L$
is initialized to the common neighbors of all vertices in $C^0$ (or
all of $V$ if $C^0$ is empty), and \texttt{CurrentDen} and
\texttt{CurrentNum} are initialized to $\den(C^0)$ and $\num(C^0)$,
respectively (or $P$ and $\widetilde{P}$ if $C^0$ is empty). Because
these sets are passed as arguments and updated incrementally as $C$
grows, redundant recomputation is avoided.

The recursion terminates when $L$ is empty
(Step~\ref{step.Mterminate}), indicating that the current clique is
maximal. In each iteration of the for-loop, a vertex $v \in L$ is
considered for addition to $C$. The sets \texttt{NewDen} and
\texttt{NewNum} are computed and the $\ell$-frequency condition is
checked. If satisfied, the incumbent is updated whenever the mortality
rate of $C \cup \{v\}$ strictly exceeds \texttt{BestObj}, and a
deeper recursive call is made with the updated sets. Regardless of
whether $v$ is added, it is removed from $L$ after its recursive call
completes. Once all recursive calls are exhausted, \texttt{BestSol}
and \texttt{BestObj} contain the optimal solution to
MMRCP~\eqref{eq.MMRCP}.

\section{A Combinatorial Branch-and-Bound
Algorithm}\label{sec.CBB}

The modified BK algorithm of Section~\ref{sec.MBK} enumerates all
$\ell$-frequent cliques, pruning only branches that violate the
$\ell$-frequency condition. Here, we strengthen this approach by
incorporating an upper-bounding scheme that enables additional
pruning, yielding a depth-first search (DFS) combinatorial
branch-and-bound (CBB) algorithm.

The upper bound is derived as follows. For an $\ell$-frequent clique
$C \supset C'$, note that $C'$ is also $\ell$-frequent, and that
$|\num(C)| \le |\num(C')|$ and $|\den(C)| \ge \ell$. Therefore,
\begin{equation}\label{eq.gammaC}
    \mu(C) \le \gamma(C') \coloneqq \frac{|\num(C')|}{\ell},
\end{equation}
which provides an upper bound on the mortality rate of any
$\ell$-frequent clique containing $C'$. This $\gamma$-upper-bound is
monotonically non-increasing: $\gamma(C'') \le \gamma(C')$ for any
$C \supset C'' \supset C'$, making it particularly effective for
pruning in DFS order, where the current clique expands monotonically
as the search descends the tree.

\begin{algorithm}[htpb]
\SetKwFunction{FDFSBB}{DfsBB}
\SetKwProg{Fn}{Function}{}{}
    \KwIn{$\langle P, \widetilde{P}, A, G=(V,E), \ell \rangle$ and
    $\ell$-frequent clique $C^0$ (possibly empty)}
    \KwOut{An $\ell$-frequent clique containing $C^0$ with maximum
    mortality rate}
        Initialize global variables $\texttt{BestSol}\gets C^0,\
        \texttt{BestObj} \gets \mu(C^0)$ \\
    \If{$C^0 = \emptyset$}{
          $v \gets \arg\max_{u \in V}\mu(u)$\\
          $\texttt{BestSol}\gets \{v\},\ \texttt{BestObj} \gets
          \mu(v)$\\
         \FDFSBB{$\emptyset, V, P, \widetilde{P}$}
    }
    \Else{
    $L \gets \bigcap\limits_{u \in C^0} N(u) \setminus C^0$\\
     \FDFSBB{$C^0, L, \den(C^0), \num(C^0)$}
    }
    \Return $\texttt{BestSol}$ and $\texttt{BestObj}$
    \Fn{\FDFSBB{$C, L, \texttt{CurrentDen},
    \texttt{CurrentNum}$}}{
        \If{$L = \emptyset$}{
            \KwRet
        }
        \For{$v \in L$}{
            $\texttt{NewDen} \gets \texttt{CurrentDen} \cap A_v$ \\
            $\texttt{NewNum} \gets \texttt{CurrentNum} \cap A_v$ \\
            \If{$|\texttt{NewDen}| \ge \ell$ and
            $\frac{|\texttt{NewNum}|}{\ell} > \texttt{BestObj}$}{
                \If{$\frac{|\texttt{NewNum}|}{|\texttt{NewDen}|} >
                \texttt{BestObj}$}{
                    $\texttt{BestObj} \gets
                    \frac{|\texttt{NewNum}|}{|\texttt{NewDen}|}$ \\
                    $\texttt{BestSol} \gets C \cup \{v\}$ \\
                }
                \FDFSBB{$C \cup \{v\}, L \cap N(v),
                \texttt{NewDen}, \texttt{NewNum}$}
            }
            $L \gets L \setminus \{v\}$
        }
        \KwRet
    }
\caption{CBB for the Maximum Mortality Rate Clique Problem}
\label{alg.CBB}
\end{algorithm}

Algorithm~\ref{alg.CBB} traverses the comorbidity graph in DFS order
and prunes branches based on $\gamma$-upper-bound violations and
failure to satisfy $\ell$-frequency. The algorithm grows the current
clique $C$ through recursive calls to \texttt{DfsBB}, each receiving
the candidate set $L$ and the denominator and numerator sets
associated with $C$. Passing these sets as arguments enables
incremental updates as $C$ is extended, avoiding recomputation from
scratch at each level. In each iteration of the for-loop, the
algorithm attempts to add a vertex $v \in L$ to $C$. Before recursing,
it verifies that $C \cup \{v\}$ is $\ell$-frequent and that
$\gamma(C \cup \{v\})$ exceeds the current incumbent objective. If
the mortality rate of the extended clique itself also improves upon
the incumbent, the incumbent solution and objective are updated
accordingly.

\section{Integer Programming Formulations}\label{sec.formulations}

This section presents two MILP formulations for MMRCP. We begin with
a natural mixed-integer nonlinear fractional program and then derive
two alternative linearizations. Throughout, a binary vector $x$
indicates which vertices are included in the clique, and a binary
vector $y$ identifies patients diagnosed with every disease in that
clique. 

Let $\widebar{G} = (V, \widebar{E})$ denote the complement of the
comorbidity graph $G = (V, E)$, and let $uv$ denote an edge $\{u,v\}$
for brevity. Suppose $G$ has $k$ connected components
$H_1,\ldots,H_k$. We describe $\calX \subseteq \{0,1\}^{|V|}$, the
set of characteristic vectors of nonempty cliques in $G$, using
binary variables $x_u$ for each vertex $u \in V$ and binary component
variables $f_i$ for $i=1,\ldots,k$:
\begin{align}
\label{eq.clqedge}  x_u + x_v &\le 1 & \forall uv \in E(\bar{H_i}),\
i=1,\ldots,k \\
\label{eq.xflink} x_u &\le f_i & \forall u \in V(H_i),\
i=1,\ldots,k \\
\label{eq.sumfto1} \sum_{i=1}^k f_i &= 1 \\
\label{eq.nonemptyclq} \sum_{u \in V} x_u &\ge 1 \\
\label{eq.xfbounds} (x,f) &\in [0,1]^{|V|} \times [0,1]^{k}.
\end{align}
Constraint~\eqref{eq.clqedge} enforces pairwise compatibility within
each component, constraint~\eqref{eq.xflink} links each selected
vertex to its component indicator, constraint~\eqref{eq.sumfto1}
ensures that selected vertices belong to exactly one component,
and constraint~\eqref{eq.nonemptyclq} enforces nonemptiness. 

MMRCP can then be formulated as:
\begin{subequations}\label{form.fracIP}
\begin{align}
  \max\ & \frac{\sum\limits_{i \in \widetilde{P}} y_i}{\sum\limits_{i
  \in P} y_i}
  \label{eq.fraclinear} \\
  \text{s.t.}\quad
  y_i &\le 1 - x_u & \forall u \in V\setminus D_i,\ i \in P
  \label{eq.xyconnectupper-disagg} \\
  y_i &\ge 1 - \sum_{u \in V\setminus D_i} x_u & \forall i \in P
  \label{eq.xyconnectlowerFP} \\
  \sum_{i \in P} y_i &\ge \ell \label{eq.denomLB} \\
  x &\in \calX \label{eq.xbin} \\
  y_i &\in \{0,1\} & \forall i \in P \label{eq.ybin}
\end{align}
\end{subequations}

The fractional objective~\eqref{eq.fraclinear} represents the
mortality rate to be maximized.
Constraint~\eqref{eq.xyconnectupper-disagg} forces $y_i = 0$ whenever
a disease $u \notin D_i$ is selected in the clique. Conversely, if
the selected clique is a subset of $D_i$,
constraint~\eqref{eq.xyconnectlowerFP} forces $y_i = 1$. Together,
these constraints ensure $y_i = 1$ if and only if patient $i$ has
every disease in the selected clique.
Constraint~\eqref{eq.denomLB} restricts attention to $\ell$-frequent
cliques.

An equivalent MILP linearizes the fractional
objective by introducing a continuous variable $z$ representing the
mortality rate and auxiliary continuous variables $w_i \in [0,1]$
representing the products $z y_i$:
\begin{subequations}\label{form.milp}
\begin{align}
  \max\ & z \label{eq.milpobj} \\
  \text{s.t.}\quad
  w_i &\le y_i & \forall i \in P \label{eq.wto0} \\
  w_i &\ge z + y_i - 1 & \forall i \in P \label{eq.wtozLB} \\
  w_i &\le z & \forall i \in P \label{eq.wtozUB} \\
  \sum_{i \in \widetilde{P}} y_i &= \sum_{i \in P} w_i
  \label{eq.modz} \\
  z &\in [0,1] \label{eq.zbound} \\
  w_i &\in [0,1] & \forall i \in P \label{eq.wbound} \\
  \nonumber&\eqref{eq.xyconnectupper-disagg}-\eqref{eq.ybin}
\end{align}
\end{subequations}

Constraints~\eqref{eq.wto0}--\eqref{eq.wtozUB} enforce $w_i = z
y_i$, and equation~\eqref{eq.modz} then captures the mortality rate
correctly in $z$.

A second linearization follows from the~\cite{charnes1962}
transformation, applicable to fractional programs with nonnegative
variables and a strictly positive denominator:
\begin{subequations}\label{form.CCIP}
\begin{align}
  \max\ & \sum_{i \in \widetilde{P}} \bar{y}_i
  \label{eq.fraclinearcc} \\
  \text{s.t.}\quad
  \bar{x}_u &\ge t - \frac{1}{|P|}(1 - x_u) & \forall u \in V
  \label{eq.xlin1CC} \\
  \bar{x}_u &\le \frac{1}{\ell} x_u & \forall u \in V
  \label{eq.xlin2CC} \\
  \bar{x}_u &\le t & \forall u \in V \label{eq.xlin3CC} \\
  \bar{y}_i &\le t - \bar{x}_u & \forall u \in V\setminus D_i,\ i
  \in P \label{eq.xyconnectupper-disaggcc} \\
  \bar{y}_i &\ge t - \sum_{u \in V\setminus D_i} \bar{x}_u &
  \forall i \in P \label{eq.xyconnectlowerFPcc} \\
  \sum_{i \in P} \bar{y}_i &= 1 \label{eq.denomCC} \\
  \bar{x}_u &\ge 0 & \forall u \in V \label{eq.barxNN} \\
  0 \le \bar{y}_i &\le \frac{1}{\ell} & \forall i \in P
  \label{eq.bary_bounds} \\
  \frac{1}{|P|} \le t &\le \frac{1}{\ell} \label{eq.t} \\
  x &\in \calX \label{eq.xbincc}
\end{align}
\end{subequations}

This reformulation introduces a continuous variable $t$ modeling the
reciprocal of the denominator of the mortality rate, i.e.,
$t = 1/\sum_{i \in P} y_i$. Since only $\ell$-frequent cliques are
of interest, constraint~\eqref{eq.t} imposes natural bounds on $t$.
The original binary variables $x \in \calX$ serve as clique incidence
vectors, and continuous variables $\bar{x}_u$ model the products
$\bar{x}_u = t x_u$ via
constraints~\eqref{eq.xlin1CC}--\eqref{eq.xlin3CC}. When $x_u = 1$,
constraints~\eqref{eq.xlin1CC} and~\eqref{eq.xlin3CC} together force
$\bar{x}_u = t$, with constraint~\eqref{eq.xlin2CC} redundant given
the upper bound in~\eqref{eq.t}. When $x_u = 0$,
constraints~\eqref{eq.xlin2CC} and~\eqref{eq.barxNN} force
$\bar{x}_u = 0$, rendering constraints~\eqref{eq.xlin1CC}
and~\eqref{eq.xlin3CC} redundant. With $\bar{x}$ correctly
represented, constraints~\eqref{eq.xyconnectupper-disaggcc}
and~\eqref{eq.xyconnectlowerFPcc} enforce $\bar{y}_i = t y_i$ for
each $i \in P$ by the same logic as
constraints~\eqref{eq.xyconnectupper-disagg}
and~\eqref{eq.xyconnectlowerFP} in
formulation~\eqref{form.fracIP}. Specifically, if $x_u = 1$ for some
$u \in V \setminus D_i$, then $\bar{x}_u = t$ and
constraint~\eqref{eq.xyconnectupper-disaggcc} forces $\bar{y}_i = 0$.
If instead $\bar{x}_u = 0$ for all $u \in V \setminus D_i$, then
$\bar{y}_i = t$. The scaling constraint~\eqref{eq.denomCC} then
ensures that objective~\eqref{eq.fraclinearcc} correctly maximizes
the mortality rate.

\section{Logic-Based Benders Reformulation}\label{sec.lbbd}

Both MILP formulations in Section~\ref{sec.formulations} linearize
the fractional objective, introducing additional variables and
constraints. While $|V|$ is moderate when diseases are represented by
ICD-coded disease categories, EHR datasets typically cover an extremely large patient
population, making $|P|$ massive.
Constraints~\eqref{eq.xyconnectupper-disagg}--\eqref{eq.denomLB}
alone, which couple the $x$ and $y$ variables and enforce
$\ell$-frequency, contribute
$2|P| + 3\sum_{i \in P}|V\setminus D_i|$ nonzero constraint
coefficients. Assuming $\min_{i \in P}|V\setminus D_i| > \epsilon|V|$,
this amounts to at least $(2+3\epsilon|V|)|P|$ nonzero entries (typically with
$\epsilon > 0.5$). These constraints appear in both MILP~\eqref{form.milp} and, in continuous form, in the Charnes--Cooper reformulation~\eqref{form.CCIP}. Their number grows with $|P|$, creating a significant computational challenge, particularly for large-scale instances when using a lazy-cut implementation of MILP~\eqref{form.milp}.

To address this scalability challenge, we introduce a
\textit{logic-based Benders reformulation} that avoids both the
$y$-variables and the linearization of the fractional objective, and
is naturally suited to decomposition~\citep{hooker2023logic}. Logic-based
Benders decomposition (LBBD) was introduced by~\cite{Hooker2003} as a
generalization of classical Benders decomposition~\citep{Benders1962},
providing a general framework for reformulating and solving challenging
combinatorial optimization problems. Whereas classical Benders
decomposition is restricted to problems whose subproblems reduce to
linear programs upon fixing a subset of variables---with Benders cuts
derived from the LP dual---LBBD extends this strategy to problems in
which the subproblem can be an arbitrary optimization
problem~\citep{Hooker2019}. The central theoretical construct enabling
this generalization is the \textit{inference dual}.

The general problem structure is as follows. Consider a maximization
problem of the form
\begin{equation}
    \max\; f(x, y) \quad \text{s.t.} \quad C(x, y),\; C^0(x),\;
    x \in D_x,\; y \in D_y,
    \label{eq.lbbd_master_orig}
\end{equation}
where $x$ denotes the complicating variables assigned by a master
problem and $y$ denotes the remaining variables resolved by a
subproblem. Fixing $x$ to $\bar{x}$ yields the subproblem
\begin{equation}
    \max\; f(\bar{x}, y) \quad \text{s.t.} \quad C(\bar{x}, y),\;
    y \in D_y.
    \label{eq.lbbd_subproblem}
\end{equation}
The inference dual for a general optimization problem
$\max\{f(x) \mid C(x),\, x \in D\}$ seeks the tightest upper bound
$v$ on the objective value that can be logically deduced from the
constraints:
\[
    \min\; v \quad \text{s.t.} \quad C(x) \underset{P}{\vdash}
    \bigl(f(x) \leq v\bigr), \quad v \in \mathbb{R},\; P \in
    \mathcal{P},
\]
where $C(x) \underset{P}{\vdash} (f(x) \leq v)$ denotes that proof
$P$, drawn from a family of proofs $\mathcal{P}$, deduces
$f(x) \leq v$ from $C(x)$~\citep{Hooker2019}. The inference dual of
subproblem~\eqref{eq.lbbd_subproblem} is accordingly
\[
    \min\; v \quad \text{s.t.} \quad C(\bar{x}, y)
    \underset{P}{\vdash} \bigl(f(\bar{x}, y) \leq v\bigr), \quad
    v \in \mathbb{R},\; P \in \mathcal{P}.
\]
When the subproblem is an LP, the inference method reduces to
nonnegative linear combinations of inequalities, and the inference
dual coincides with the classical LP dual~\citep{Hooker2019}.
Let $v^*$ denote the optimal value of the subproblem for a given
$\bar{x}$, and let $P^*$ be the proof establishing
$f(\bar{x}, y) \leq v^*$. The key insight of LBBD is that the same
proof $P^*$ may establish a valid bound $B_{\bar{x}}(x)$ on the
subproblem optimal value for other values of $x$, satisfying
$B_{\bar{x}}(\bar{x}) = v^*$. Two types of Benders cuts arise
depending on the subproblem outcome~\citep{Hooker2019}. When the
subproblem at iteration $i$ (with master assignment $x^i$) is feasible,
the bound $B_{x^i}(x)$ is added to the master problem as an
\textit{optimality cut}
\[
    z \leq B_{x^i}(x),
\]
which tightens the master problem objective by enforcing that the
subproblem value cannot exceed $B_{x^i}(x)$ for any assignment
of $x$. When the subproblem is infeasible, a \textit{feasibility cut}
$N_{x^i}(x)$ is derived instead, which excludes the current master
solution $x^i$ from the feasible region. At iteration $k$, the
master problem incorporating both cut types is
\[
    z_k = \max\; z \quad \text{s.t.} \quad C^0(x);\;
    z \leq B_{x^i}(x),\; i \in \mathcal{F};\;
    N_{x^i}(x),\; i \in \mathcal{I};\;
    x \in D_x,
\]
where $\mathcal{F}$ and $\mathcal{I}$ denote the index sets of
iterations yielding feasible and infeasible subproblems,
respectively~\citep{Hooker2019}. The optimal value $z_k$ constitutes
a valid upper bound on problem~\eqref{eq.lbbd_master_orig}, while the
best subproblem value found thus far provides a lower bound. The
algorithm terminates when these bounds coincide.

Because LBBD cuts must be designed specifically for each problem class
rather than being automatically derived from an LP dual, their
construction requires problem-specific insight~\citep{Hooker2019}.
This necessity can also be regarded as an advantage: carefully crafted
cuts can be substantially tighter than generic LP-dual cuts, yielding
convergence in considerably fewer iterations. Furthermore, the
framework naturally accommodates heterogeneous solver combinations~\citep{Hooker2019}. LBBD also forms the
theoretical basis for the combinatorial Benders cuts of
\citet{Codato2006}, which exploit the LBBD framework to reformulate
and solve MILP formulations containing big-M constraints used to model
conditional implications. A comprehensive study of LBBD theory,
inference duality, and applications is provided in
\citet{Hooker2019} and \citet{hooker2023logic}.

Notable applications of LBBD span a broad range of combinatorial
optimization domains. In scheduling, it has been applied to avionics
scheduling with partial assignment acceleration~\citep{KARLSSON2022}
and preemptive flexible job-shop
scheduling~\citep{JUVIN2023106156}. In healthcare operations, it has
been used for stochastic distributed operating room scheduling via
binary decision diagrams~\citep{Guo2020} and for chronic outpatient
scheduling via answer set programming~\citep{CAPPANERA2023}. Further
applications include multi-manned assembly line balancing with
symmetry-breaking cuts~\citep{MICHELS2019796} and two-dimensional bin
packing using area-based models and no-good
cuts~\citep{Cote2020}.

A fundamental challenge in applying LBBD to MMRCP is the non-monotone
behavior of the mortality rate as a function of the selected disease
clique (Theorem~\ref{thm.mu-nonmonotonic}). This necessitates the
use of \textit{simple no-good} optimality cuts, which are known to be
extremely weak~\citep{hooker2023logic}. Strengthening these cuts and
embedding the reformulation within a decomposition branch-and-cut
(DBC) algorithm is therefore a central focus of this work. The
feasibility cuts, by contrast, are cover inequalities based on minimal
non-$\ell$-frequent cliques, which are considerably stronger.

We now present the reformulation, which uses the $x$-variables to
indicate the selected clique and a scalar variable $z$ to capture the
mortality rate. We first introduce additional notation. For binary
vectors $x, \widetilde{x} \in \{0,1\}^{|V|}$, the Hamming distance
$\calH(x, \widetilde{x})$ counts the number of positions in which $x$
and $\widetilde{x}$ differ:
\[
\mathcal{H}(x, \widetilde{x}) = \sum_{u \in J_0(\widetilde{x})} x_u
+ \sum_{u \in J_1(\widetilde{x})} (1-x_u),
\]
where $J_0(\widetilde{x}) \coloneqq \{u \in V : \widetilde{x}_u = 0\}$
and $J_1(\widetilde{x}) \coloneqq \{u \in V : \widetilde{x}_u = 1\}$.
For notational convenience, we overload the mortality rate and
$\gamma$-upper-bound functions as follows: $\mu(x) \coloneqq
\mu(J_1(x))$, $\mu(uv) \coloneqq \mu(\{u,v\})$, $\gamma(uv)
\coloneqq \gamma(\{u,v\})$, $\mu(u) \coloneqq \mu(\{u\})$, and
$\gamma(u) \coloneqq \gamma(\{u\})$.

The logic-based Benders reformulation of MMRCP is:
\begin{subequations}\label{form.LBBF}
\begin{align}
  \max\ & z \\
  \text{s.t.}\quad
  z &\le \mu(\widetilde x) + (1 - \mu(\widetilde x))
  \mathcal{H}(x,\widetilde{x}) & \forall \widetilde{x} \in \calX_o
  \label{eq.LBBoptcut}\\
  \sum_{u \in J_1(\widetilde{x})} x_u &\le |J_1(\widetilde{x})| - 1
  & \forall \widetilde{x} \in \calX_f \label{eq.LBBfeascut}\\
  x &\in \calX,
\end{align}
\end{subequations}
where
\begin{align*}
    \calX_o &\coloneqq \{\widetilde{x} \in \calX :
    J_1(\widetilde{x}) \text{ is an $\ell$-frequent clique}\}, \\
    \calX_f &\coloneqq \{\widetilde{x} \in \calX :
    J_1(\widetilde{x}) \text{ is a minimal non-$\ell$-frequent
    clique}\}.
\end{align*}

Constraint~\eqref{eq.LBBoptcut} is the simple no-good optimality cut:
when $x = \widetilde{x}$, it forces $z \le \mu(x)$, correctly
computing the objective for that feasible solution, and is satisfied
for all $z \in [0,1]$ whenever $x \ne \widetilde{x}$.
Constraint~\eqref{eq.LBBfeascut} is a cover inequality that serves as
a feasibility cut eliminating non-$\ell$-frequent cliques. It exploits
the fact that every superset of a non-$\ell$-frequent clique is itself
non-$\ell$-frequent: for cliques $C \subset C'$ in $G$,
$|\den(C)| \le \ell - 1$ implies $|\den(C')| \le \ell - 1$. The cut
forces at least one vertex of any minimal non-$\ell$-frequent clique
to be excluded from the solution.

To use this reformulation effectively within a DBC algorithm, several
implementation details must be addressed, including the master problem
at the root node, cuts added to strengthen the relaxation, and other
enhancements discussed next.

\subsection{Root Relaxation and Delayed Constraint
Generation}\label{sec.gamma-lbbd}

The optimality cuts~\eqref{eq.LBBoptcut} and feasibility
cuts~\eqref{eq.LBBfeascut} in reformulation~\eqref{form.LBBF} can be
exponentially many in the worst
case~\citep{MoonMoser65}, as they correspond to nonempty cliques in
the comorbidity graph. As is standard in Benders decomposition, these
cuts are generated in a delayed fashion: a cut is added only when the
optimal solution to a relaxation violates it. Since the initial
relaxation without these cuts is an integer program, it is solved via
branch-and-cut (BC), and violated cuts are added at BC tree nodes upon
encountering integral solutions. This scheme constitutes the DBC
algorithm.

Given the weakness of the simple no-good optimality cuts, the strength
of the root node relaxation is particularly important. We therefore
develop valid inequalities based on the $\gamma$-upper-bounds for
single vertices $v$ and edge endpoints $uv$, intended for use at the
root node. The root node relaxation,
formulation~\eqref{form.GamLBBD}, combines $\gamma$-upper-bound-based
valid inequalities when they are non-trivial (i.e., less than one)
with the corresponding special cases of the simple no-good
cuts~\eqref{eq.LBBoptcut} otherwise:
\begin{subequations}\label{form.GamLBBD}
\begin{align}
  \max\ & z \\
  \text{s.t.}\quad
  \label{eq.single} z &\le \mu(v) + (1 - \mu(v)) \mathcal{H}(x,
  e_v) & \forall v \in V : \gamma_v \ge 1 \\
  z &\le \gamma_v x_v + (1 - x_v) & \forall v \in V :
  \gamma_v < 1 \\
  z &\le \mu(uv) + (1 - \mu(uv)) \mathcal{H}(x, e_u + e_v) &
  \forall uv \in E : \gamma(uv) \ge 1 \\
  \label{eq.pair} z &\le \gamma(uv)(x_u + x_v - 1) +
  (2 - x_u - x_v) & \forall uv \in E : \gamma(uv) < 1 \\
  \label{eq.zmuLB} z &\ge \mu_{LB} \\
  x &\in \calX,
\end{align}
\end{subequations}
where $\mu_{LB} \coloneqq \max_{u \in V} \mu(u)$ and $e_u$ denotes
the $|V|$-dimensional unit vector with a one in position $u$.

\begin{algorithm}[htbp]
\KwIn{$\widetilde{x} \in \calX$ and $\widetilde{z} \in [0,1]$}
\KwOut{Violated optimality or feasibility cut, if one exists}
\If{$J_1(\widetilde{x})$ is not $\ell$-frequent}{
    $\widetilde{C} \gets$ \texttt{Minimalize}{($J_1(\widetilde{x})$)}\\
    \Return violated feasibility cut:
        $\sum\limits_{u \in \widetilde{C}}x_u \le |\widetilde{C}| - 1$
}
\ElseIf{$\widetilde{z} > \mu(\widetilde{x})$}{
    \Return violated optimality cut:
        $z \le \mu(\widetilde{x}) + (1-\mu(\widetilde{x})) \times
        \mathcal{H}(x, \widetilde{x})$
}
\Else{\Return no violated cuts found}
\caption{Logic-Based Benders Decomposition: Separation Procedure}
\label{alg.LBBDgamma}
\end{algorithm}

The DBC algorithm begins by solving formulation~\eqref{form.GamLBBD}
at the root node. At each BC tree node, if an integral solution to the
LP relaxation is obtained, the separation procedure in
Algorithm~\ref{alg.LBBDgamma} checks feasibility with respect to
reformulation~\eqref{form.LBBF}. 
As established in
Remark~\ref{prop.f1-clique-guarantee} below, $J_1(\widetilde{x})$ is
guaranteed to be a clique. A feasibility cut is added if it is not
$\ell$-frequent; otherwise, a violated optimality cut is added if one
exists. Violated cuts can be identified in polynomial time. If a violated cut is found, the node LP relaxation
is re-solved; otherwise, the node is pruned by feasibility. The
separation procedure is implemented via the solver's callback
functionality to add lazy cuts. Algorithm~\ref{alg.minimalize}
describes how a minimal non-$\ell$-frequent clique is identified to
generate cover-based feasibility cuts. The resulting algorithm (Algorithm~\ref{alg.LBBDgamma}) is referred to henceforth as the
\textit{pure LBBD} algorithm. 

\begin{remark}\label{prop.f1-clique-guarantee}
Let $\calX_{LP} = \{x : (x,f) \text{ satisfies }
\eqref{eq.clqedge}-\eqref{eq.xfbounds}\}$ denote the LP relaxation
of $\calX$. If $x^* \in \calX_{LP}$, then
$C \coloneqq \{v \in V : x^*_v > \tfrac{1}{2}\}$ is a clique in $G$.
To see this, suppose for contradiction that $C$ contains non-adjacent
vertices $u$ and $v$, with $u \in V(H_i)$ and $v \in V(H_j)$.
Constraints~\eqref{eq.xflink} imply $f_i^* \ge x_u^* > \tfrac{1}{2}$
and $f_j^* \ge x_v^* > \tfrac{1}{2}$, which violates either
constraint~\eqref{eq.clqedge} (if $i = j$) or
constraint~\eqref{eq.sumfto1} (if $i \ne j$).
\end{remark}

\begin{algorithm}[htbp]
\KwIn{Clique $\widetilde{C}$ that is not $\ell$-frequent}
\KwOut{Minimal non-$\ell$-frequent clique}
    \SetKwFunction{FMinimal}{Minimalize}
    \SetKwProg{Fn}{Function}{}{}
    \Fn{\FMinimal{$\widetilde{C}$}}{
        \Repeat{\texttt{stopflag} = \texttt{true}}{
            \texttt{stopflag} = \texttt{true}\\
            \For{$u \in \widetilde{C}$}{
                \If{$|\den(\widetilde{C}\setminus \{u\})| < \ell$}{
                    $\widetilde{C} \gets \widetilde{C} \setminus
                    \{u\}$\\
                    \texttt{stopflag} = \texttt{false}\\
                }
            }
        }
        \KwRet $\widetilde{C}$
    }
\caption{Clique Minimalization}
\label{alg.minimalize}
\end{algorithm}

\subsection{Hybrid LBBD Approach}\label{sec.hybridLBBD}

This section describes hybrid strategies that use the CBB algorithm
from Section~\ref{sec.CBB} to generate strong valid inequalities for
incorporation into the LBBD DBC algorithm. The theoretical basis for
this hybrid approach is given by Theorem~\ref{thm.hybridF1cut}.

\begin{theorem}[CBB cut]\label{thm.hybridF1cut}
Let $\widetilde{C}$ be an $\ell$-frequent clique, and let $\bar{C}$
be an $\ell$-frequent clique containing $\widetilde{C}$ with maximum
mortality rate. If $x$ corresponds to an arbitrary $\ell$-frequent
clique in $G$ and $z = \mu(x)$, then $(x,z)$ satisfies:
\begin{equation} \label{eq.hybrid-cutF1}
z \le \mu(\bar{C}) + (1-\mu(\bar{C})) \sum_{u \in \widetilde{C}}
(1 - x_u).
\end{equation}
\end{theorem}

\begin{proof}
If $\widetilde{C} \subseteq J_1(x)$, then $\mu(x) \le \mu(\bar{C})$
by definition of $\bar{C}$, since $x$ corresponds to an
$\ell$-frequent clique containing $\widetilde{C}$. Otherwise,
$\widetilde{C} \setminus J_1(x) \ne \emptyset$, so
$\sum_{u \in \widetilde{C}} (1 - x_u) \ge 1$ and the right-hand side
of~\eqref{eq.hybrid-cutF1} is at least one, making the inequality
trivially satisfied.
\end{proof}

Algorithm~\ref{alg.hybridLBBD} describes the separation procedure
incorporating the CBB cut. At BC tree nodes where an integral solution
is found, the LBBD algorithm invokes the CBB cut whenever the clique
size meets or exceeds a user-specified threshold $\tau$.

\begin{algorithm}[htbp]
\KwIn{$\widetilde{x} \in \calX$, $\widetilde{z} \in [0,1]$, and
positive integer $\tau$}
\KwOut{Violated optimality, CBB, or feasibility cut, if detected}
\If{$J_1(\widetilde{x})$ is not $\ell$-frequent}{
    $\widetilde{C} \gets$ \texttt{Minimalize}{($J_1(\widetilde{x})$)}\\
    \Return violated feasibility cut:
        $\sum\limits_{u \in \widetilde{C}}x_u \le |\widetilde{C}| - 1$
}
\ElseIf{$\widetilde{z} > \mu(\widetilde{x})$ and
$|J_1(\widetilde{x})| < \tau$}{
    \Return violated optimality cut:
        $z \le \mu(\widetilde{x}) + (1-\mu(\widetilde{x})) \times
        \mathcal{H}(x, \widetilde{x})$
}
\ElseIf{$\widetilde{z} > \mu(\widetilde{x})$ and
$|J_1(\widetilde{x})| \ge \tau$}{
    $\bar{C}, \bar{\mu} \gets$ Output of Algorithm~\ref{alg.CBB}
    with $C^0 = J_1(\widetilde{x})$\\
    \Return violated optimality cut:
        $z \le \mu(\widetilde{x}) + (1-\mu(\widetilde{x})) \times
        \mathcal{H}(x, \widetilde{x})$\\
    and CBB cut: $z \le \bar{\mu} + (1-\bar{\mu})
    \sum\limits_{u \in J_1(\widetilde{x})} (1 - x_u)$
}
\Else{\Return no violated cuts found}
\caption{Hybrid Logic-Based Benders Decomposition: Separation
Procedure}
\label{alg.hybridLBBD}
\end{algorithm}

The hybrid approach benefits from the CBB algorithm while retaining
the expressiveness of the integer programming formulation. Its
implementation in Algorithm~\ref{alg.hybridLBBD} rests on two
observations. First, the CBB call---which is the most computationally
expensive step---is triggered only when the current integral solution
contains at least $\tau$ vertices, ensuring it is invoked only when a
sufficiently structured $\ell$-frequent clique is already available.
In this case, the CBB algorithm operates on the subgraph induced by
$J_1(\widetilde{x}) \cup \bigcap_{u \in J_1(\widetilde{x})} N(u)$,
a considerably smaller graph than $G$, which keeps the cost of each
CBB call manageable. Second, CBB is invoked only at integral solutions
that already violate the simple no-good cut, so the number of CBB
calls is implicitly controlled by the branch-and-cut process.

The CBB algorithm is also used to strengthen the root node relaxation
by generating CBB cuts~\eqref{eq.hybrid-cutF1} for the $k$ vertices
with the highest individual mortality rates, for a user-specified $k$,
and adding them to formulation~\eqref{form.GamLBBD}. The best
mortality rate clique found across these $k$ calls also improves the
lower bound $\mu_{LB}$ in constraint~\eqref{eq.zmuLB}. The resulting
algorithm is referred to henceforth as the \textit{hybrid LBBD}
algorithm.

The methods developed in this chapter are evaluated on real-world EHR
data. Chapter~\ref{sec.DATAPREP} describes the preparation of the testbed,
including the extraction and preprocessing of data from the MIMIC-IV
database, the mapping of diagnosis codes to higher-level disease
categories, and the construction of the comorbidity graph on which all
experiments are conducted.